\section{Ograniczenia, falsyfikacja i dalsza praca}

Wynik tej pracy należy czytać jako kontrfaktyczny eksperyment obliczeniowy, a nie punktową prognozę makroekonomiczną dla Polski. \amor{} bada ścieżki alternatywne względem wariantu bez BDP i pokazuje mechanizm: przy deficytowym finansowaniu transfer może zwiększać presję automatyzacyjną tylko do momentu, w którym kanał płacowo-rezerwowy i fiskalno-finansowy zaczynają ograniczać zdolność firm do finansowania inwestycji. Dlatego centralna interpretacja opiera się na różnicach względem ścieżki bez BDP, nie na absolutnych poziomach terminalnych.

\subsection{Diagnostyka wariantu bez BDP}

Model startuje ze skalibrowanego punktu Polski 2026-04-30, ale wartości z pierwszego miesiąca obejmują już przepływy i księgowania wykonane przez model. Zewnętrzne punkty odniesienia pochodzą z raportu inflacyjnego NBP, uchwały RPP z 4 marca 2026 r., tablic GUS, strategii długu Ministerstwa Finansów oraz Article IV MFW dla Polski \parencite{nbp2026inflationReport,nbp2026rppResolution,gus2026unemploymentMarch,gus2026gdp2025Preliminary,mf2025debtStrategy,imf2026polandArticleIV}.

\begin{table}[H]
\centering
\scriptsize
\caption{Diagnostyka wariantu bez BDP: punkt odniesienia zewnętrzny, pierwszy wynik miesięczny i ścieżka po 60 miesiącach}
\label{tab:baseline-plausibility}
\begin{tabularx}{\textwidth}{@{}L{0.20\textwidth}L{0.20\textwidth}rrX@{}}
\toprule
Zmienna & Punkt odniesienia zewnętrzny & Model, m1 & Bez BDP, m60 & Interpretacja \\
\midrule
Inflacja CPI & 3,0\% w III 2026 / projekcja blisko celu & 3,33\% & 2,90\% & wejście i koniec ścieżki pozostają blisko celu inflacyjnego \\
Stopa referencyjna NBP & 3,75\% od 05.03.2026 & 3,55\% & 3,08\% & m1 odzwierciedla regułę reakcji modelu, nie literalny punkt decyzji RPP \\
Bezrobocie & 6,1\% rejestrowane, koniec III 2026 & 6,09\% & 7,40\% & start bliski kotwicy; wariant bez BDP generuje pogorszenie rynku pracy \\
Dług publiczny/PKB & ok. 53--66\%, zależnie od definicji PDP/EDP & 47,09\% & 69,34\% & ścieżka bez BDP generuje wzrost długu; poziom terminalny nie jest prognozą punktową \\
Deficyt/PKB & ok. 6--7\% w prognozach 2025--2026 & 13,59\% & 7,69\% & m1 obejmuje startowe przepływy wykonania modelu; porównania BDP opierają się na różnicach względem wariantu bez BDP \\
Rachunek bieżący/PKB & około -1\% PKB w projekcji MFW 2026 & -5,23\% & -9,75\% & poziom bezwzględny nie jest prognozą; interpretowane są zmiany względem tej samej ścieżki bez BDP \\
Inwestycje prywatne/PKB & GUS: stopa inwestycji ogółem 17,0\% PKB w 2025; prywatny odpowiednik jest węższy & 11,44\% & 9,90\% & zmienna modelowa jest węższa od GFCF ogółem; używana głównie jako kontrfaktyczna delta \\
\bottomrule
\end{tabularx}
\end{table}


Tabela~\ref{tab:baseline-plausibility} pokazuje, że wariant bez BDP jest wiarygodnym punktem startowym dla części zmiennych krótkookresowych, zwłaszcza inflacji, stopy NBP i bezrobocia rejestrowanego. Jednocześnie swobodna ścieżka modelu generuje własny dryf fiskalno-zewnętrzny: dług i deficyt pozostają wysokie, a rachunek bieżący pogarsza się bardziej, niż sugerowałyby zewnętrzne projekcje. To ograniczenie modelu, nie ukryty wynik BDP; dlatego rachunek bieżący i inwestycje należy czytać jako odchylenia od tej samej ścieżki bez BDP.

\subsection{Status twierdzeń empirycznych}

Tabela~\ref{tab:claim-status} porządkuje wynik falsyfikacyjnie: model odrzuca hipotezę monotoniczną i wspiera wersję warunkową, zależną od zdolności firm do finansowania reakcji technologicznej.

\begin{table}[htbp]
\centering
\scriptsize
\caption{Status głównych twierdzeń po eksperymencie BDP}
\label{tab:claim-status}
\begin{tabularx}{\textwidth}{@{}>{\raggedright\arraybackslash}p{0.34\textwidth}>{\raggedright\arraybackslash}p{0.18\textwidth}>{\raggedright\arraybackslash}X@{}}
\toprule
\textbf{Twierdzenie} & \textbf{Status} & \textbf{Uzasadnienie w wynikach} \\
\midrule
BDP zwiększa automatyzację monotonicznie: im wyższy transfer, tym większa adopcja AI/hybrid. &
Odrzucone &
Adopcja rośnie do okolic BDP 2000~PLN, po czym spada przy BDP 2500~PLN i BDP 3000~PLN. \\
\addlinespace
Deficytowo finansowany BDP może działać jako warunkowy katalizator automatyzacji. &
Wspierane &
W wariancie centralnym istnieje przedział dodatniego efektu względem wariantu bez BDP; efekt zanika po przekroczeniu progu absorpcji. \\
\addlinespace
Kanał płacy rezerwowej samodzielnie tłumaczy wynik. &
Niewspierane &
Analiza \(\lambda\) pokazuje, że kanał płacy rezerwowej jest ważny, ale sam nie wyjaśnia lokalnego maksimum adopcji. \\
\addlinespace
Nieliniowość wyniku wynika z załamania gospodarstw domowych. &
Niewspierane &
Panel gospodarstw domowych pokazuje stabilne metryki dystrybucyjne, brak bankructw gospodarstw i poprawę płynności wraz z BDP. \\
\addlinespace
Sam kanał fiskalno-finansowy wystarcza do ograniczenia nakładów inwestycyjnych. &
Niewspierane &
Dla \(\lambda = 0\) deficyt i dług rosną, ale adopcja AI/hybrid oraz inwestycje pozostają blisko wariantu bez BDP. \\
\addlinespace
Interakcja kanału płacowo-rezerwowego i fiskalno-finansowego ogranicza zdolność firm do finansowania nakładów inwestycyjnych. &
Wspierane &
Przy wyższych \(\lambda\) presja automatyzacyjna rośnie równolegle z pogorszeniem finansowania, spadkiem inwestycji i załamaniem adopcji po lokalnym maksimum. \\
\addlinespace
Eksperyment dowodzi, że realna Polska po BDP osiągnęłaby dokładnie takie wartości makroekonomiczne. &
Odrzucone &
Model służy analizie kontrfaktycznej ścieżki i mechanizmów transmisji, a nie punktowej prognozie. \\
\bottomrule
\end{tabularx}
\end{table}

\subsection{Ograniczenia}

Pierwsze ograniczenie dotyczy samego scenariusza BDP. Model obejmuje 100 000 agentów reprezentujących gospodarstwa domowe, a nie pełną populację 38 mln osób, więc transfer jest skalowanym impulsem fiskalno-dochodowym, nie administracyjną mikrosymulacją per osoba, per dorosły i per dziecko. Wszystkie główne wnioski dotyczą wariantu deficytowo finansowanego; VAT, PIT, składka od dochodu, cięcia wydatków albo podatki majątkowe mogłyby zmienić kanały transmisji \parencite{owsiak2017}.

Drugie ograniczenie dotyczy mostu MFW--\amor{}. Raport MFW analizuje redystrybucję i dobrobyt, a dla Polski opiera się na danych z 2013 r. Ta praca wykorzystuje MFW jako punkt odniesienia metodologiczny, nie jako scenariusz SFC-ABM ani alternatywną estymację efektów redystrybucyjnych. Wartości MFW pomagają zakotwiczyć skalę transferu, ale porównanie ma charakter interpretacyjny, nie rachunkowy.

Trzecie i czwarte ograniczenie dotyczą zachowań oraz technologii. Analiza odporności obejmuje jeden wymiar behawioralny, czyli \(\lambda\), i nie testuje alternatywnych reguł podaży pracy ani pełnej heterogeniczności demograficznej gospodarstw. Miara bankructw gospodarstw jest kontrolą pomocniczą, nie modelem prawnym upadłości konsumenckiej. Model mierzy adopcję AI/hybrid, ale nie rozróżnia robotyzacji produkcji, automatyzacji zaplecza administracyjnego, systemów predykcyjnych i generatywnej AI; wynik należy więc czytać jako analizę kanałów makro-finansowych wpływających na zdolność automatyzacji, nie jako prognozę konkretnych technologii.

Piąte i szóste ograniczenie dotyczą sektora publicznego oraz skali eksperymentu. Model zawiera reguły fiskalne i monetarne, ale nie opisuje procesu politycznego, wiarygodności rządu, reakcji wyborców ani korekty programu BDP w trakcie jego działania. Wariant centralny obejmuje 10 ziaren losowania i horyzont 60 miesięcy. To wystarcza do raportowania średnich, odchyleń i parowanych przedziałów ufności dla głównych kontrfaktów, ale nie zastępuje programu obejmującego wiele wariantów finansowania, parametrów behawioralnych i szoków zewnętrznych.

\subsection{Dalsza praca}

Najważniejsze dalsze kroki to porównanie sposobów finansowania BDP, pełniejsza heterogeniczność gospodarstw domowych i podaży pracy, rozbicie miary AI/hybrid na konkretne klasy technologii oraz replikacja scenariusza w \amor{}. Celem takiego programu byłoby rozdzielenie efektu transferu od efektu finansowania i utrzymanie modelu jako narzędzia porównywania założeń, a nie zamkniętego źródła ostatecznych odpowiedzi.
